% ===========================================================================
% AUTOSCIENCE HOUSE PAPER TEMPLATE — two-column, math-dense workshop paper.
% Matches the reference look: full-width title + abstract, 2-column justified
% body, Computer Modern, numbered sections, booktabs tables, numbered display
% equations, dense inline math, single- and full-width figures.
%
% Compiles offline with pdflatex (no exotic packages). When a specific venue
% .sty is selected for a run, ADAPT to it instead of this template — keep the
% same structure and math density.
% ===========================================================================
\documentclass[10pt,twocolumn]{article}
\usepackage[letterpaper,margin=0.85in,columnsep=0.28in]{geometry}
\usepackage{amsmath,amssymb,amsfonts}
\usepackage{booktabs}
\usepackage{graphicx}
\usepackage[font=small,labelfont=bf]{caption}
\usepackage{microtype}
\usepackage{url}

\title{\bf <Paper Title>}
\author{<Author> \\ <Affiliation> \\ \texttt{<email or ``email not provided''>}}
\date{<Month Year>}

\begin{document}

% --- full-width title + abstract (spans both columns) ----------------------
\twocolumn[
  \begin{@twocolumnfalse}
    \maketitle
    \begin{abstract}
      \noindent
      % Abstract opens with the one-sentence thesis, then the scoped empirical
      % claim. Keep it honest and specific; quote the headline numbers exactly
      % as repro.sh produces them.
      <one-sentence thesis>. <scoped empirical claim with exact numbers>.
    \end{abstract}
    \vspace{1.0em}
  \end{@twocolumnfalse}
]

\section{Introduction}
% Narrow scoped contribution. Include an explicit "what we did NOT find in a
% literature scan on <date>" sentence, and a one-line differentiation of this
% paper's axis vs sibling papers. Use inline math freely, e.g. the independent
% variable is the configuration $c = (\text{optimizer}, \lambda, \dots)$.

\section{Method}
% Precise and reproducible: equations, grid/config, exact hyperparameters.
% Use numbered display equations for every defined quantity, e.g.
\begin{equation}
  D(d,\sigma) \;=\; \frac{1}{T\binom{N}{2}} \sum_{t=1}^{T} \sum_{i<j}
  \mathbf{1}\!\left[ S_{i,t}(d,\sigma) \neq S_{j,t}(d,\sigma) \right],
  \label{eq:metric}
\end{equation}
% and define every symbol immediately after.

\section{Results}
% Figures + booktabs tables; honest reporting of censored/partial outcomes.
% Single-column figure:
% \begin{figure}[t]\centering\includegraphics[width=\linewidth]{figures/fig1}
%   \caption{...}\label{fig:1}\end{figure}
% booktabs table:
\begin{table}[t]
  \centering
  \caption{Headline results. Every value is produced by \texttt{repro.sh}.}
  \label{tab:headline}
  \small
  \begin{tabular}{lcc}
    \toprule
    Setting & Metric & 95\% CI \\
    \midrule
    <row> & <value> & <[lo, hi]> \\
    \bottomrule
  \end{tabular}
\end{table}
% For wide tables/figures spanning both columns use table*/figure*.

\section{Checks and Robustness}
% Controls pass, seed-independence, threshold sensitivity. State the control
% values (e.g. $D(0)=0$) explicitly.

\section{Discussion and Limitations}
% State what is NOT claimed. Failure modes reported as findings, not hidden.

\section{Conclusion}
% One paragraph: the scoped takeaway.

\section*{Reproducibility statement}
% Seeds, per-cell counts, one-command repro (\texttt{bash repro.sh}), and the
% concrete artifact/release list.

\begin{thebibliography}{9}
\bibitem{ref1} <Author>. <Title>. <Venue>, <year>.
\end{thebibliography}

\end{document}
